%% 自定义命令
% fun[x范围][y范围][路径名称,或其它算子参数][表達式]
\NewDocumentCommand\fun{  O{xmin:xmax} O{ymin:ymax}  O{name} O{} m }{%
\addplot+ [raw gnuplot, empty line = jump, no markers, name path = #3]%
        gnuplot {%
        set xrange [#1];
        set yrange [#2];
        set contour base;
        set cntrparam levels discrete 0.003;
        unset surface;
        % set view map;
        % set isosamples 50;
        set samples 50;
        splot #5;
    }  #4;
}
% 
\NewDocumentCommand\calxy{r()}{%
    \footnotesize
    \pgfplotspointgetcoordinates{(#1)}%
    (\pgfmathprintnumber[/pgf/number format/precision=2]{\pgfkeysvalueof{/data point/x}},%
    \pgfmathprintnumber[/pgf/number format/precision=2]{\pgfkeysvalueof{/data point/y}})
}

% 
\NewDocumentCommand\calx{r()}{%
\pgfplotspointgetcoordinates{(#1)}%
\pgfmathprintnumber[/pgf/number format/precision=2]%
{\pgfkeysvalueof{/data point/x}}%
}
% 
\NewDocumentCommand\caly{r()}{%
\footnotesize\pgfplotspointgetcoordinates{(#1)}%
\pgfmathprintnumber[/pgf/number format/precision=2]%
{\pgfkeysvalueof{/data point/y}}%
}



\NewDocumentCommand\mathexp{m m}{%
node[pos = #1, pin = #2]{}
}
% 
\def\xmin{\pgfkeysvalueof{/pgfplots/xmin}}
\def\xmax{\pgfkeysvalueof{/pgfplots/xmax}}
\def\ymin{\pgfkeysvalueof{/pgfplots/ymin}}
\def\ymax{\pgfkeysvalueof{/pgfplots/ymax}}